\chapter{Θεωρητικό Υπόβαθρο Εδάφους και Υγρασίας}
\label{ch:theory}

\section{Φυσικές Ιδιότητες του Εδάφους και Μέθοδοι Μέτρησης Υγρασίας}

\subsection{Εισαγωγή}

Στο κεφάλαιο αυτό παρουσιάζονται οι βασικές έννοιες του εδάφους και
περιγράφονται οι κύριες φυσικές και χημικές του ιδιότητες. Επιπλέον,
παρουσιάζονται οι μορφές ύπαρξης του νερού εντός του εδάφους, καθώς και
οι βασικές μέθοδοι που χρησιμοποιούνται για τον υπολογισμό της
περιεκτικότητάς του σε νερό.

\subsection{Ορισμός και Χαρακτηριστικά Εδάφους}

Το έδαφος αποτελεί το ανώτερο στρώμα του γήινου φλοιού και απαρτίζεται
από στερεά υλικά (οργανική και ανόργανη ύλη), νερό και αέρια. Πρόκειται
για ένα δυναμικά μεταβαλλόμενο σύστημα που ορίζεται από πολλές φυσικές
παραμέτρους όπως η υφή, η θερμοκρασία, η αλατότητα και το pH \cite{ref01}.

\paragraph{Υφή}

Μία από τις σημαντικότερες ιδιότητες είναι η υφή του εδάφους. Ανάλογα με
τη σύσταση των ανόργανων υλικών (άμμος, ιλύς και άργυλος) τα εδάφη
ταξινομούνται σε μία από τις παρακάτω κατηγορίες υφής: αμμώδες (sandy),
ιλυώδες (silty), αργιλώδες (clayey), πηλώδες (loamy), αμμοπηλώδες (sandy
loam), πηλοαμμώδες (loamy sand), ιλυοπηλώδες (silty loam), αργιλοπηλώδες
(clay loam), ιλυοαργιλοπηλώδες (silty clay loam), αμμοαργιλοπηλώδες
(sandy clay loam), ιλυοαργιλώδες (silty clay), αμμοαργιλώδες (sandy
clay). \cite{ref02}

Η υφή του εδάφους θεωρείται σχετικά μόνιμο χαρακτηριστικό και επηρεάζει
σημαντικά τις ιδιότητες του, όπως τη συγκράτηση του νερού, τον αερισμό
αλλά και τη διείσδυση των ριζών. Για το λόγο αυτό είναι ιδιαίτερα
χρήσιμος ο χαρακτηρισμός του εδάφους στη γεωργία, καθώς διαδραματίζει
καθοριστικό ρόλο στην ανάπτυξη των φυτών.

\begin{figure}[H]
  \centering
  \includegraphics[width=0.56\textwidth]{soil-texture-triangle.png}
  \caption{Διάγραμμα υφής του USDA που απεικονίζει τις δώδεκα κατηγορίες υφής σύμφωνα με τα ποσοστά αργίλου, ιλύος και άμμου. \cite{ref01}}
  \label{fig:soil-texture-triangle}
\end{figure}

\paragraph{Τιμή pH}

Μία ακόμη χημική ιδιότητα είναι το pH, το οποίο εκφράζει πόσο όξινο ή
αλκαλικό είναι το έδαφος. Η γνώση του pH θεωρείται ιδιαίτερα σημαντική,
καθώς επηρεάζει τη διαθεσιμότητα των θρεπτικών συστατικών προς τα φυτά,
τη δραστηριότητα των μικροοργανισμών αλλά και τον καθορισμό της
καταλληλότητας του εδάφους για διαφορετικές καλλιέργειες. Η διαλυτότητα
και κατά συνέπεια η διαθεσιμότητα των διαφόρων χημικών στοιχείων
μεταβάλλεται ανάλογα με την τιμή του pH, καθορίζοντας έτσι το εύρος
τιμών στο οποίο μπορεί να ευδοκιμήσει κάθε καλλιέργεια, ανάλογα με τις
θρεπτικές της απαιτήσεις. \cite{ref07}

\paragraph{Αλατότητα}

Το έδαφος εκτός από τη στερεά μορφή του περιλαμβάνει και την υγρή φάση,
το επονομαζόμενο εδαφικό διάλυμα, στο οποίο συναντώνται διαλυμένα άλατα
με τη μορφή ιόντων. Τα κύρια στοιχεία είναι το άζωτο, ο φώσφορος και το
κάλιο. Τα ιόντα αυτά αποτελούν την άμεσα διαθέσιμη μορφή θρεπτικών
στοιχείων επηρεάζοντας τη γονιμότητά του \cite{ref02} . Εκτός από τις φυσικές
του ιδιότητες, τα ιόντα επηρεάζουν και τις ηλεκτρικές του ιδιότητες,
καθώς η ηλεκτρική αγωγιμότητα του εδάφους χρησιμοποιείται συνήθως ως
ένδειξη του πλήθους των διαλυμένων αλάτων. Το μέγεθος αυτό βοηθάει στην
εκτίμηση της ποιότητας του εδάφους, καθώς η παρουσία διαλυμένων αλάτων
είναι απαραίτητη για την ανάπτυξη των φυτών, ενώ σε υψηλές συγκεντρώσεις
μπορεί να μειώσουν τη διαθεσιμότητα του νερού \cite{ref09}.

\paragraph{Πορώδες}

Η παρουσία νερού στο έδαφος είναι κρίσιμη, καθώς είναι ο μηχανισμός
μεταφοράς των θρεπτικών ουσιών. Ένα μέρος αποθηκεύεται στους πόρους του
εδάφους κι ένα μέρος αποστραγγίζεται. Η διαθεσιμότητά του προς τα φυτά
εξαρτάται από το πόσο ισχυρά συγκρατείται από το εδαφικό υλικό. Στη
διαδικασία αυτή σημαντικό ρόλο διαδραματίζει η δομή του εδάφους, και
ειδικότερα το πόσο πορώδες είναι. Σε συμπιεσμένα εδάφη, περιορίζεται η
διήθηση του νερού και εμποδίζεται η ανάπτυξη του ριζικού συστήματος. Σε
ένα σωστά δομημένο έδαφος διευκολύνεται η κατάλληλη διήθηση και
αποστράγγιση του νερού, καθώς και η κυκλοφορία του αέρα. \cite{ref08}

\begin{figure}[H]
  \centering
  \includegraphics[width=0.55\textwidth]{porosity-infiltration.png}
  \caption{Διήθηση νερού σε έδαφος με μειωμένο και αυξημένο πορώδες αντίστοιχα.}
  \label{fig:porosity-infiltration}
\end{figure}

\subsection{Καταστάσεις νερού στο Έδαφος}

Η ύπαρξη νερού, αποτελεί μεταβλητή αλλά απαραίτητη ιδιότητα του εδάφους.
Στο εδαφικό διάλυμα η κατάσταση του νερού υπάρχει σε δύο βασικές μορφές,
το \textbf{Ελεύθερο Νερό} (Free Water) και το \textbf{Δεσμευμένο Νερό}
(Bound Water) οι οποίες απεικονίζονται στο \figref{fig:free-bound-water}. Όταν το νερό
βρεθεί κοντά στην επιφάνεια των κόκκων, ένα μέρος απορροφάται και
συγκρατείται ισχυρά από δυνάμεις προσρόφησης, αποτελώντας το δεσμευμένο
νερό. Σε αντίθεση με το ελεύθερο, το δεσμευμένο έχει μειωμένη
κινητικότητα και δε μπορεί εύκολα να προσανατολιστεί όταν εφαρμοστεί
ηλεκτρικό πεδίο \cite{ref03,ref04}.

\begin{figure}[H]
  \centering
  \includegraphics[width=0.44\textwidth]{free-bound-water.png}
  \caption{Διαχωρισμός του νερού στο έδαφος σε Ελεύθερο Νερό και Δεσμευμένο νερό. \cite{ref04}}
  \label{fig:free-bound-water}
\end{figure}

Το νερό, όπως και κάθε υλικό, χαρακτηρίζεται από τη \textbf{διηλεκτρική
σταθερά}, ένα μέγεθος ιδιαίτερα χρήσιμο για τον προσδιορισμό της
ποσότητας νερού στο έδαφος. Η διηλεκτρική σταθερά εξαρτάται από την
ικανότητά του υλικού να πολώνεται υπό την επίδραση ηλεκτρικού πεδίου .
Ως αποτέλεσμα, το ελεύθερο νερό έχει μεγαλύτερη διηλεκτρική σταθερά σε
σχέση με το δεσμευμένο νερό, τα μόρια του οποίου ευθυγραμμίζονται
δυσκολότερα με το πεδίο \cite{ref04}. Η ποσότητα του νερού που θα δεσμευθεί
εξαρτάται άμεσα και από τον τύπο του εδάφους, με ενδεικτικό παράδειγμα
τα αργιλώδη εδάφη να συγκρατούν περισσότερο νερό συγκριτικά με τα αμμώδη
λόγω της μεγαλύτερης επιφάνειας των κόκκων\cite{ref03}.

\subsection{Έκφραση της Περιεκτικότητας του Εδάφους σε Νερό}

Η περιεκτικότητα του εδάφους σε νερό εκφράζεται με δύο βασικά μεγέθη:
την \textbf{βαρυμετρική} περιεκτικότητα σε νερό (Gravimetric Water
Content ή GWC) και την \textbf{ογκομετρική} περιεκτικότητα σε νερό
(Volumetric Water Content ή VWC).

Η βαρυμετρική περιεκτικότητα σε νερό είναι η μάζα του νερού ανά μονάδα
μάζας εδάφους, δηλαδή γραμμάρια νερού ανά γραμμάρια εδάφους και
υπολογίζεται σύμφωνα με τον παρακάτω τύπο

\[w = \frac{M_{moist}\,(g) - M_{dry}\,(g)}{M_{dry}\,(g)}\]

Για τον υπολογισμό του ζυγίζεται αρχικά ένα δείγμα υγρού εδάφους και στη
συνέχεια εφόσον ξηρανθεί και έχει εξατμιστεί όλη η ποσότητα του νερού,
πραγματοποιείται και πάλι η ζύγιση του δείγματος. Αποτελεί μία από τις
κύριες μεθόδους για τον υπολογισμό της υγρασίας καθώς η ποσότητα του
νερού υπολογίζεται απευθείας με τη μέτρηση των μαζών. \cite{ref05}

Η \textbf{ογκομετρική περιεκτικότητα} σε νερό ορίζεται ως ο όγκος του
νερού προς το συνολικό όγκο του εδαφικού δείγματος, όπως φαίνεται από
τον παρακάτω τύπο

\[\theta = \frac{V_{water}\,\left( m^{3} \right)}{V_{soil}\,\left( m^{3} \right)}\]

Η ογκομετρική περιεκτικότητα προκύπτει είτε έμμεσα από τη βαρυμετρική,
είτε χρησιμοποιείται κάποιος αισθητήρας για τον υπολογισμό της. Οι
αισθητήρες αυτοί αξιοποιούν τις διηλεκτρικές ιδιότητες του νερού, αφού
το νερό λόγω της υψηλής διηλεκτρικής σταθεράς του θα αλληλεπιδράσει
έντονα με το ηλεκτρομαγνητικό πεδίο που εφαρμόζεται στο δείγμα από τον
αισθητήρα. Τελικά, σύμφωνα με την απόκριση των μορίων του νερού στο
πεδίο γίνεται η εκτίμηση της ογκομετρικής υγρασίας. \cite{ref06}

Η βαρυμετρική και η ογκομετρική περιεκτικότητα εκφράζουν το ίδιο φυσικό
μέγεθος, με το w να βασίζεται στη μάζα ενώ το θ στον όγκο. Συσχετίζοντας
τα δύο μεγέθη, η ογκομετρική περιεκτικότητα μπορεί να εκφραστεί ως το
γινόμενο της βαρυμετρικής με τη φαινόμενη ξηρή πυκνότητα του εδάφους
ρ\textsubscript{b} (dry bulk density)

\[\theta = w\,\rho_{b}\]

όπου ρ\textsubscript{b} ορίζεται ως ο λόγος της μάζας του ξηρού εδάφους
(m\textsubscript{d}) προς το συνολικό όγκο του εδαφικού δείγματος (
V\textsubscript{soil} ) \cite{ref05,ref06}

\[\rho_{b} = \frac{m_{d}\,(g)}{V_{soil}\,\left( cm^{3} \right)}\]

\section{Ηλεκτρομαγνητική Συμπεριφορά Εδάφους}

\subsection{Εισαγωγή}

Το χώμα αποτελεί ένα ετερογενές και πορώδες υλικό. Η ηλεκτρομαγνητική
συμπεριφορά ενός μη-σιδηρομαγνητικού υλικού, όπως και του εδάφους,
μπορεί να παρομοιαστεί ως ένας πυκνωτής με απώλειες (lossy capacitor).
Το σήμα που περνάει μέσα από το χώμα θα υποστεί ιοντικές απώλειες οι
οποίες οφείλονται σε διαλυμένα άλατα (Na,Cl) τα οποία κινούνται ελεύθερα
υπό την επίδραση του πεδίου. Έτσι, εμφανίζεται η αγωγιμότητα που
οφείλεται στο ρεύμα μέσω της πραγματικής μεταφοράς ηλεκτρικού φορτίου
\cite{ref10}.

Παράλληλα, τα μόρια του νερού λόγω της διπολικότητάς τους, θα
προσπαθήσουν να ευθυγραμμιστούν με το εναλλασσόμενο πεδίο. Αυτό δίνει
στο μέσο έναν διηλεκτρικό χαρακτήρα, ο οποίος θα εμφανίσει το ρεύμα
μετατόπισης \cite{ref10}.


Στο παρόν κεφάλαιο, θα αναλυθούν οι δύο συνιστώσες που καθορίζουν την
ηλεκτρομαγνητική απόκριση του εδάφους και είναι χρήσιμες για τον
προσδιορισμό σημαντικών φυσικών παραμέτρων: η αγώγιμη και τη διηλεκτρική
συνιστώσα .

\subsection{Αγωγιμότητα και Ειδική Αντίσταση -- Μέθοδος Wenner}

Ως ηλεκτρική αντίσταση ενός αγωγού ορίζεται η δυσκολία που παρουσιάζει
στη διέλευση του ηλεκτρικού ρεύματος, κι εκφράζεται από το νόμο του Οhm
ως:

\[R = \frac{V}{I},\quad(\Omega)\]

όπου V είναι η διαφορά δυναμικού στα άκρα του αγωγού σε Volt και I η
ένταση του ρεύματος σε Ampere. Επειδή ωστόσο η αντίσταση εξαρτάται άμεσα
από τη γεωμετρία του αγωγού, εκφράζεται επίσης από τη σχέση:

\[R = \rho\frac{L}{A},\quad(\Omega)\]

όπου ρ η ειδική αντίσταση του αγωγού σε Ω*m, L το μήκος του αγωγού και Α
το εμβαδό της διατομής του. Το αντίστροφο της αντίστασης ονομάζεται
ηλεκτρική αγωγιμότητα και συμβολίζεται ως \cite{ref14}:

\[G = \frac{1}{R},\quad(Siemens)\]

ενώ το αντίστροφο της ειδικής αντίστασης ονομάζεται ειδική ηλεκτρική
αγωγιμότητα και συμβολίζεται ως \cite{ref13}:

\[\sigma = \frac{1}{\rho},\quad(Siemens/m)\]

Η ειδική ηλεκτρική αγωγιμότητα χρησιμοποιείται συχνά ως δείκτης για τη
συγκέντρωση των διαλυμένων αλάτων στο έδαφος (αλατότητα). Η αγωγιμότητα
του εδάφους συμβολίζεται ως EC και εκφράζεται σε Siemens/m.

Μία διαδεδομένη μέθοδος μέτρησης της αλατότητας του εδάφους στηρίζεται
στον υπολογισμό της ειδικής αντίστασης με τη μέθοδο Wenner. Κατά τη
μέθοδο αυτή, χρησιμοποιούνται τέσσερα ηλεκτρόδια που τοποθετούνται στο
έδαφος σε ευθεία γραμμή σε ίσες συνήθως μεταξύ τους αποστάσεις
(\figref{fig:wenner-array}). Στο ζευγάρι των εξωτερικών ηλεκτροδίων γίνεται η έγχυση
ρεύματος, ενώ το εσωτερικό ζευγάρι χρησιμοποιείται για τη μέτρηση της
διαφοράς δυναμικού \cite{ref15}.

\begin{figure}[H]
  \centering
  \includegraphics[width=0.68\textwidth]{wenner-array.png}
  \caption{Η διάταξη Wenner με τα τέσσερα ηλεκτρόδια A, M, N και Β σε απόσταση a, τις ισοδυναμικές γραμμές και τη ροή του ρεύματος εξ αιτίας του ηλεκτρομαγνητικού πεδίου. Το ρεύμα ρέει από τα εξωτερικά ηλεκτρόδια, ενώ η πτώση τάσης μετράται από τα εσωτερικά. \cite{ref14}}
  \label{fig:wenner-array}
\end{figure}

Επειδή η πραγματική ειδική αντίσταση του μέσου θα διαφέρει από τη
μέτρηση, χρησιμοποιείται ο συμβολισμός της φαινόμενης ειδικής αντίστασης
ρ\textsubscript{α}. Για απόσταση a και βάθος c, ο τύπος που δίνει το
ρ\textsubscript{α} είναι \cite{ref11}:

\[\rho_{\alpha} = \frac{4\pi\alpha R}{1 + \frac{2\alpha}{\sqrt{\alpha^{2} + 4c^{2}}} - \frac{\alpha}{\sqrt{\alpha^{2} + c^{2}}}}\]

Επομένως,

\[EC_{\alpha} = \sigma_{\alpha} = \frac{1}{\rho_{\alpha}} = \frac{1 + \frac{2\alpha}{\sqrt{\alpha^{2} + 4c^{2}}} - \frac{\alpha}{\sqrt{\alpha^{2} + c^{2}}}}{4\pi\alpha R}\]

Εάν θεωρήσουμε το βάθος των ηλεκτροδίων πολύ μικρότερο από τη μεταξύ
τους απόσταση, ο παραπάνω τύπος μπορεί να απλοποιηθεί ως εξής \cite{ref12}:

\[\rho_{\alpha} = 2\pi\alpha R\]

\subsection{Θεωρία των Διηλεκτρικών}

\textbf{Διηλεκτρικά} ονομάζονται τα υλικά που είναι μη αγώγιμα στο
ηλεκτρικό ρεύμα, αλλά χαρακτηρίζονται από την ικανότητα των δεσμευμένων
ηλεκτρικών φορτίων τους να \textbf{πολώνονται} όταν εκτίθενται σε
ηλεκτρικό πεδίο. Πρακτικά, η ροή του ρεύματος είναι μικρή ή σχεδόν
μηδενική ενώ η ενέργεια αποθηκεύεται στο ηλεκτρικό πεδίο που
δημιουργείται εντός του διηλεκτρικού λόγω της πόλωσης. \textbf{Ηλεκτρική
Πόλωση (P)} είναι το μέγεθος που εκφράζει το βαθμό διαχωρισμού των
θετικών και αρνητικών φορτίων ενός υλικού όταν αυτό εκτεθεί σε ηλεκτρικό
πεδίο \cite{ref17}.

Η ικανότητα αυτή ενός συστήματος να αποθηκεύει ηλεκτρικό φορτίο
εκφράζεται από τη \textbf{χωρητικότητα}. Πιο συγκεκριμένα, πρόκειται για
το μονόμετρο φυσικό μέγεθος που συναντάται συνήθως στους πυνκωτές και
ορίζεται ως το πηλίκο του ηλεκτρικού φορτίου Q του πυκνωτή ως προς την
τάση V μεταξύ των οπλισμών του πυκνωτή \cite{ref15},

\[C = \frac{Q\,(Coulomb)}{V\,(Volt)},\quad Farad\]

Ο παραπάνω τύπος εφαρμόζεται για κάθε μορφής πυκνωτή, ωστόσο η
χωρητικότητα εξαρτάται από το σχήμα, τις διαστάσεις, την απόσταση των
οπλισμών, καθώς και το διηλεκτρικό που παρεμβάλλεται μεταξύ των
οπλισμών. Γι'αυτό το λόγο, για έναν επίπεδο πυκνωτή με αέρα ή κενό ως
διηλεκτρικό, η χωρητικότητα δίνεται από τη σχέση \cite{ref28}

\[C = \frac{\varepsilon_{0}\, A}{d},\quad F\quad\quad\text{(10)}\]

όπου $\varepsilon_{0}=8{,}85\cdot 10^{-12}\,\mathrm{C^2/(N\,m^2)}$ η \textbf{απόλυτη
διηλεκτρική σταθερά} του κενού, Α το εμβαδόν του οπλισμού και d η
απόσταση των οπλισμών. Η \textbf{απόλυτη διηλεκτρική σταθερά ή απόλυτη
διηλεκτρική διαπερατότητα} αποτελεί ένα μέγεθος το οποίο εκφράζει την
ικανότητα ενός διηλεκτρικού υλικού να πολώνεται στο εσωτερικό του όταν
εφαρμοστεί ηλεκτρικό πεδίο. Ειδικότερα, ορίζεται ως ο λόγος της
ηλεκτρικής μετατόπισης \(\overrightarrow{D}\)
\(\left( \frac{C}{m^{2}} \right)\) προς το ηλεκτρικό πεδίο
\(\overrightarrow{E}\left( \frac{V}{m} \right)\) στο εσωτερικό του
\cite{ref16}

\[\varepsilon = \frac{\overrightarrow{D}}{\overrightarrow{E}},\quad F/m\]

Εάν μεταξύ των οπλισμών υπάρχει διηλεκτρικό υλικό, τότε η εξίσωση (10)
μετατρέπεται σε \cite{ref15,ref28}

\[C = \frac{\varepsilon_{0}\,\varepsilon_{r}\, A}{d},\quad F\]

όπου ε\textsubscript{r} η \textbf{σχετική διηλεκτρική σταθερά ή σχετική
διηλεκτρική διαπερατότητα} του υλικού.

Η εισαγωγή του διηλεκτρικού έχει ως αποτέλεσμα την αύξηση της
χωρητικότητας συγκριτικά με την ύπαρξη κενού. Αυτό συμβαίνει λόγω της
πόλωσης των φορτίων του υλικού, τα οποία δημιουργούν εσωτερικό πεδίο
αντίθετης φοράς από το εφαρμοζόμενο πεδίο. Το μέτρο της απόκρισης ενός
υλικού στο εφαρμοζόμενο πεδίο χαρακτηρίζεται από τη \textbf{σχετική
διηλεκτρική σταθερά} (relative dielectric constant) \textbf{ή}
\textbf{σχετική διηλεκτρική διαπερατότητα} (relative dielectric
permittivity) και λαμβάνει μία δεδομένη τιμή που εξαρτάται από το
εκάστοτε υλικό \cite{ref15}.

Η σχετική διηλεκτρική σταθερά συνδέεται με την απόλυτη διηλεκτρική
σταθερά μέσω της σχέσης \cite{ref15}

\[\varepsilon_{r} = \frac{\varepsilon}{\varepsilon_{0}}\]

\begin{figure}[H]
  \centering
  \includegraphics[width=0.78\textwidth]{dielectric-polarization.png}
  \caption{Απεικόνιση του μη πολωμένου υλικού απούσια εξωτερικού πεδίου στα αριστερά. Στα δεξιά απεικονίζεται η πόλωση των διπόλων του υλικού κατά την εφαρμογή ηλεκτρικού πεδίου, με την πόλωση να δίνεται ως \(\overrightarrow{P} = \left( \varepsilon-\varepsilon_{0} \right) \cdot \overrightarrow{E}\) . \cite{ref16}}
  \label{fig:dielectric-polarization}
\end{figure}

\subsection{Μη ιδανικά Διηλεκτρικά}

Σε μη ιδανικά διηλεκτρικά, όπως και το έδαφος, ταυτόχρονα με την πόλωση
παρουσιάζονται και αγώγιμες απώλειες οι οποίες τελικά μετατρέπονται σε
θερμότητα. Η συνολική αυτή απόκριση του διηλεκτρικού σε εναλλασσόμενο
πεδίο μπορεί να περιγραφεί από την \textbf{απόλυτη} \textbf{ηλεκτρική
μιγαδική διαπερατότητα ε* (absolute permittivity)} \cite{ref18}:

\[\varepsilon^{*}(\omega) = \varepsilon'(\omega) + j\varepsilon_{eff}''(\omega)\]

και συνδέεται με τη σχετική διηλεκτρική διαπερατότητα με τη σχέση:

\[\varepsilon^{*}(\omega) = \varepsilon_{0}\,\varepsilon_{r}^{*}(\omega)\]

όπου \(\varepsilon_{0}\) η ηλεκτρική διαπερατότητα του κενού. Από τα
παραπάνω προκύπτει η μιγαδική \textbf{σχετική ηλεκτρική διαπερατότητα} η
οποία αποτελεί αδιάστατο μέγεθος και περιγράφει την αλληλεπίδραση ενός
υλικού σε σχέση με το πεδίο \cite{ref18,ref19}:

\[\varepsilon_{r}(\omega) = \varepsilon_{r}'(\omega) + j\varepsilon_{r,eff}''(\omega)\]

To εr'(ω) αποτελεί το πραγματικό μέρος της μιγαδικής διαπερατότητας και
εκφράζει το ποσό της ενέργειας που έχει αποθηκευτεί στο υλικό από το
εξωτερικό πεδίο. Το ε'\,'reff(ω) αποτελεί το φανταστικό μέρος και
σχετίζεται με τις διηλεκτρικές και θερμικές απώλειες. Στα μη ιδανικά
διηλεκτρικά όπου συνυπάρχουν οι απώλειες λόγω πόλωσης και ηλεκτιρκής
αγωγιμότητας, εκφράζονται ως \cite{ref18,ref19}:

\[\varepsilon_{r,eff}''(\omega) = \varepsilon_{r,d}''(\omega) + \frac{\sigma}{\omega\,\varepsilon_{0}}\]

όπου σ(S/m) η DC αγωγιμότητα και ε\textsubscript{rd}'\,' το φανταστικό
μέρος της σχετικής διαπερατότητας που εκφράζει τις απώλειες λόγω των
μηχανισμών πόλωσης. Επειδή, ωστόσο, στους περισσότερους τύπους εδαφών το
ε\textsubscript{rd}'\,' μπορεί να θεωρηθεί σχετικά μικρό, η μέτρηση του
φανταστικού μέρους χρησιμοποιείται συχνά για την εκτίμηση της
αγωγιμότητας σε συχνότητες στην περιοχή των MHz \cite{ref19}. Επομένως,
λαμβάνοντας υπόψη ότι τα μεγέθη αυτά εξαρτώνται από τη συχνότητα, χάριν
απλούστευσης χρησιμοποιείται η έκφραση:

\[\varepsilon_{r}^{*} = \varepsilon_{r}' + j\varepsilon_{r}''\]

\[\varepsilon_{r}'' = \frac{\sigma}{\omega\,\varepsilon_{0}}\]

\subsection{Το Υγρό Έδαφος ως Διηλεκτρικό}

Το μόριο του νερού αποτελείται από ένα άτομο οξυγόνου (Ο) το οποίο φέρει
μερικό αρνητικό φορτίο (δ-) και δύο άτομα υδρογόνου (Η) τα οποία φέρουν
μερικό θετικό φορτίο. Λόγω της γωνιακής του δομής και της
ηλεκτραρνητικότερης φύσης του οξυγόνου, δημιουργείται ένα μόνιμο
ηλεκτρικό δίπολο όπως φαίνεται στο \figref{fig:water-dipole} \cite{ref19}

\begin{figure}[H]
  \centering
  \includegraphics[width=0.62\textwidth]{water-dipole.png}
  \caption{Το διπολικό μόριο του νερού με τις μερικές φορτίσεις (δ+ και δ-). \cite{ref19}}
  \label{fig:water-dipole}
\end{figure}

Με την εφαρμογή εναλλασσόμενου ηλεκτρικού πεδίου, το δίπολο
περιστρέφεται ώστε να ευθυγραμμιστεί με το πεδίο. Στην περιοχή
συχνοτήτων MHz έως και χαμηλά GHz, τα μόρια του νερού ακολουθούν τις
ηλεκτρομαγνητικές ταλαντώσεις του πεδίου, ενώ τα σωματίδια του εδάφους
παραμένουν κατά κύριο λόγο ακλόνητα. Η συμπεριφορά αυτή δικαιολογεί τη
σημαντικά υψηλότερη σχετική διηλεκτρική σταθερά του νερού (περίπου 80
στους 25\textsuperscript{ο}) σε σύγκριση με αυτή των στερεών συστατικών
του εδάφους (περίπου 4-5 στους 25\textsuperscript{ο}) \cite{ref19}
(\figref{fig:soil-component-permittivity}).

\begin{figure}[H]
  \centering
  \includegraphics[width=0.72\textwidth]{soil-component-permittivity.png}
  \caption{Σχετική διηλεκτρική σταθερά των κύριων συστατικών του εδάφους και του νερού. \cite{ref21}}
  \label{fig:soil-component-permittivity}
\end{figure}

Οι περισσότεροι σύγχρονοι αισθητήρες βασίζονται σε αυτή την ιδιότητα των
μορίων του νερού προκειμένου να υπολογίσουν την περιεκτικότητά του σε
ένα δείγμα εδάφους. Συνεπώς, το μέγεθος που χρησιμοποιείται για την
εκτίμηση της υγρασίας στο έδαφος είναι το πραγματικό μέρος της σχετικής
διαπερατότητας ε'\textsubscript{r}. Η προσθήκη νερού στο έδαφος αυξάνει
σημαντικά τη συνολική ικανότητα του μέσου να αποθηκεύει ενέργεια μέσω
των μορίων του νερού, φαινόμενο που αντικατοπτρίζεται από την πραγματική
συνιστώσα ε'r.

Αντίθετα, το ε'\,' αποτελεί τη φανταστική συνιστώσα και όπως αναφέρθηκε
συνδέεται με τις απώλειες ενέργειας του μέσου. Συνδέεται στενά με το σ
μέσω της εξίσωσης~(2.1) και αντικατοπτρίζει την αγωγιμότητα λόγω της
ροής ελεύθερων ιόντων. Οι δύο συνιστώσες εξαρτώνται από τη συχνότητα.
Στο \figref{fig:permittivity-vs-frequency} φαίνεται η συμπεριφορά των εr' και εr΄΄ ενός μη ιδανικού
διηλεκτρικού που περιέχει αγώγιμα μονοπάτια, όπως το έδαφος.

\begin{figure}[H]
  \centering
  \includegraphics[width=0.73\textwidth]{permittivity-vs-frequency.png}
  \caption{Μεταβολή της πραγματικής και φανταστικής συνιστώσας της σχετικής διαπερατότητας και περιοχές εμφάνισης των μηχανισμών πόλωσης σε συνάρτηση με τη γωνιακή συχνότητα. \cite{ref18}}
  \label{fig:permittivity-vs-frequency}
\end{figure}

Το έδαφος αποτελείται από στερεά και υγρά σώματα, αέρα και άλατα, με
αποτέλεσμα η ροή του ρεύματος και η πόλωση των μορίων να μην είναι ίδια
σε όλο το φάσμα συχνοτήτων. Έτσι, στη διαμόρφωση της ηλεκτρικής
μιγαδικής διαπερατότητας συμβάλλει ένα σύνολο από μηχανισμούς πολώσεων.
Στο γράφημα παρουσιάζονται τέσσερις κύριοι μηχανισμοί \cite{ref20}:

\begin{itemize}
\item
  H διεπιφανειακή πόλωση (Interfacial Polarization) που σχετίζεται με τη
  συσσώρευση\\
  φορτίων σε διεπιφάνειες περιοχών λόγω της ανομοιογένειας του μέσου
\item
  Η διπολική πόλωση (Orientational Polarization) λόγω περιστροφής
  μόνιμων δίπολων όπως το νερό
\item
  Η ιοντική πόλωση (ionic polarization) λόγω ταλάντωσης των ιόντων του
  υλικού
\item
  Η ηλεκτρονική πόλωση (electronic polarization) που συμβαίνει σε
  ουδέτερα άτομα όταν
\item
  το ηλεκτρικό πεδίο μετατοπίζει των πυρήνα σε σχέση με τα ηλεκτρόνια
  που τον περιβάλλουν
\end{itemize}

Στις χαμηλές συχνότητες η απόκριση του εδάφους επηρεάζεται έντονα από
την ιοντική αγωγιμότητα που οφείλεται στην πραγματική μεταφορά ιόντων
κυρίως μέσα στο νερό των πόρων. Η αγώγιμη αυτή απώλεια, αποτυπώνεται στο
ε'\,'reff. Όπως αποτυπώνεται και από τον όρο σ/(ω*εο), η επίδραση της
αγωγιμότητας είναι μεγαλύτερη στις χαμηλότερες συχνότητες και μειώνεται
όσο αυξάνει η συχνότητα. Για το λόγο αυτό, η χαμηλή περιοχή συχνοτήτων
είναι ιδιαίτερα ευαίσθητη στην αγώγιμη συμπεριφορά του εδάφους και
χρησιμοποιείται συνήθως για την εκτίμηση της αλατότητας. Ωστόσο, στη
χαμηλή περιοχή συχνοτήτων συνυπάρχει και η διεπιφανειακή πόλωση, ενώ με
την αύξηση της συχνότητας ο μηχανισμός δεν προλαβαίνει να ακολουθήσει το
πεδίο κι εξασθενεί \cite{ref18,ref20}.

Από την άλλη, το ε' εμφανίζει υψηλές τιμές σε χαμηλές συχνότητες, καθώς
τα κινητά φορτία έχουν το χρόνο να μετακινηθούν και να και να
συσσωρευτούν στις διαφορετικές διεπιφάνειες του ανομοιογενές μέσου
δημιουργώντας χωρητικά φαινόμενα (διεπιφανειακή πόλωση). Αντίθετα, σε
υψηλότερες συχνότητες ο μηχανισμός αυτός δε δύναται να ακολουθήσει το
πεδίο κι εξασθενεί, μειώνοντας το ε'. Έτσι, αναδεικνύεται και
απομονώνεται καλύτερα η ικανότητα του υλικού να αποθηκεύει ενέργεια
κυρίως μέσω της πόλωσης των μορίων του νερού και διευκολύνει την
εκτίμηση του εr' (διπολική πόλωση).

Σε ακόμη υψηλότερες συχνότητες, έξω από το συνηθισμένο εύρος μέτρησης
υγρασίας, λαμβάνουν χώρα άλλοι μηχανισμοί όπως η ιοντική και ηλεκτρονική
πόλωση \cite{ref18,ref20}.

\begin{figure}[H]
  \centering
  \includegraphics[width=0.84\textwidth]{polarization-mechanisms.png}
  \caption{Σχηματική απεικόνιση μηχανισμών πόλωσης σε διηλεκτρικό υλικό σε μικροσκοπικό επίπεδο. \cite{ref19}}
  \label{fig:polarization-mechanisms}
\end{figure}

\subsection{Μέθοδοι Μέτρησης της Υγρασίας Εδάφους με τη χρήση Αισθητήρων}

Η ανάπτυξη των αισθητήρων υγρασίας βασίζεται κυρίως στην εφαρμογή
ηλεκτρικού πεδίου και στην αλληλεπίδραση των μορίων του νερού με αυτό.
Κάποιες από τις κύριες τεχνολογίες που αξιοποιούν την αρχή αυτή
παρουσιάζονται παρακάτω.

Οι αισθητήρες \textbf{TDR} (Time Domain Reflectometry) λειτουργούν
σύμφωνα με τον υπολογισμό του χρόνου που χρειάζεται ένας
ηλεκτρομαγνητικός παλμός για να διαδοθεί κατά μήκος δύο ή τριών συνήθως
μεταλλικών ράβδων τοποθετημένων στο δείγμα, οι οποίες λειτουργούν ως
γραμμή μεταφοράς. Η ταχύτητα διάδοσης του παλμού εξαρτάται από τη
διηλεκτρική διαπερατότητα του δείγματος, η οποία επηρεάζεται άμεσα από
την περιεκτικότητά του σε νερό. Τελικά, από τη διηλεκτρική διαπερατότητα
εκτιμάται η ογκομετρική περιεκτικότητα μέσω εμπειρικών σχέσεων
βαθμονόμησης. Στο \figref{fig:tdr-setup} φαίνεται η πειραματική διάταξη και στο
\figref{fig:tdr-waveform} φαίνεται η διάδοση του αρχικού παλμού και του ανακλόμενου
σήματος. \cite{ref23}

\begin{figure}[H]
  \centering
  \includegraphics[width=0.58\textwidth]{tdr-setup.png}
  \caption{Πειραματική διάταξη για τη μέτρηση υγρασίας στο πεδίο του χρόνου με TDR. \cite{ref23}}
  \label{fig:tdr-setup}
\end{figure}

\begin{figure}[H]
  \centering
  \includegraphics[width=0.58\textwidth]{tdr-waveform.png}
  \caption{Διάδοση του αρχικού παλμού και του ανακλώμενου σήματος κατά τη μέτρηση με αισθητήρα TDR. \cite{ref23}}
  \label{fig:tdr-waveform}
\end{figure}

Μία ακόμη μέθοδος είναι η \textbf{FDR} (Frequency Domain Reflectometry)
η οποία βασίζεται στη μέτρηση του συντελεστή ανάκλασης ενός
ηλεκτρομαγνητικού κύματος στο πεδίο της συχνότητας. Με τη χρήση ενός
Αναλυτή Διανυσματικών Δικτύων (Vector Network Analyzer), το δείγμα
διεγείρεται από ένα ηλεκτρομαγνητικό σήμα μεταβαλλόμενης συχνότητας και
καταγράφεται το πλάτος και η φάση (\figref{fig:fdr-vna-setup}). \cite{ref23}

\begin{figure}[H]
  \centering
  \includegraphics[width=0.58\textwidth]{fdr-vna-setup.png}
  \caption{Πειραματική διάταξη για τη μέτρηση υγρασίας στο πεδίο της συχνότητας (FDR) με Αναλυτή Διανυσματικών Δικτύων. \cite{ref23}}
  \label{fig:fdr-vna-setup}
\end{figure}

Το μετρούμενο μέγεθος είναι ο Σύνθετος Συντελεστής Ανάκλασης Γ, ο οποίος
περιγράφει το πόσο ισχυρά ανακλάται το σήμα σε κάθε συχνότητα και
περιέχει πληροφορία τόσο για το μέγεθος όσο και για τη φάση της
ανάκλασης. Από την κατάλληλη επεξεργασία του συντελεστή ανάκλασης μπορεί
να εξαχθεί πληροφορία σχετικά με τις ηλεκρομαγνητικές ιδιότητες του
υλικού, όπως η σχετική ηλεκτρική διαπερατότητα που όπως αναφέρθηκε
πρόκειται για έναν ισχυρό δείκτη εκτίμησης για την περιεκτικότητα του
νερού \cite{ref23}.

Παρά την υψηλή ακρίβεια που προσφέρουν οι παραπάνω μέθοδοι, τόσο το
κόστος υλοποίησης όσο και η πολυπλοκότητα του συστήματος είναι ιδιαίτερα
αυξημένα. Μία απλουστευμένη και χαμηλότερου κόστους προσέγγιση που
χρησιμοποιείται συχνά βασίζεται στη \textbf{Σύνθετη ή Μιγαδική
Αντίσταση} όπως εφαρμόζεται στα άρθρα \cite{ref24,ref25,ref26,ref27}, μέθοδος που
θα ακολουθηθεί και στην παρούσα εργασία.

Πιο συγκεκριμένα, ένα εναλλασσόμενο σήμα δεδομένης συχνότητας
εφαρμόζεται στο υπό εξέταση δείγμα και μετρώνται η πτώση τάσης και η
διαφορά φάσης μεταξύ τάσης και ρεύματος. Στη συνέχεια, από τις μετρήσεις
αυτές υπολογίζεται η σύνθετη αντίσταση του δείγματος συσχετίζεται με τη
μιγαδική σχετική ηλεκτρική διαπερατότητα του υλικού. Για τη μετάβαση από
το ηλεκτρικό αυτό μέγεθος σε ποσοτική έκφραση της περιεκτικότητας του
δείγματος σε νερό, είναι απαραίτητη η εξαγωγή ενός ηλεκτρικού ισοδύναμου
μοντέλου του δείγματος. Μέσω αυτού, υπολογίζεται η σχετική ηλεκτρική
διαπερατότητα του δείγματος, η οποία εξαρτάται άμεσα από την
περιεκτικότητά του σε νερό. Το θεωρητικό υπόβαθρο της τεχνικής αυτής θα
αναλυθεί στις επόμενες ενότητες.

\begin{figure}[H]
  \centering
  \includegraphics[width=0.72\textwidth]{impedance-measurement-two-electrode.png}
  \caption{Αναπαράσταση της μέτρησης της μιγαδικής εμπέδησης εδάφους με διάταξη δύο ηλεκτροδίων και εφαρμογή εναλλασσόμενου σήματος.}
  \label{fig:impedance-two-electrode}
\end{figure}

Όπως είναι φανερό, ανεξάρτητα από τη μέθοδο που ακολουθείται, ο τελικός
στόχος αφορά τον υπολογισμό της σχετικής ηλεκτρικής διαπερατότητας του
δείγματος. Τελικά, η παράμετρος αυτή συσχετίζεται με την κατάλληλη
εμπειρική σχέση βαθμονόμησης, με την πλέον διαδεδομένη στη βιβλιογραφία
να αποτελεί η εξίσωση \textbf{Topp}, και περιγράφεται από τη συνάρτηση
\cite{ref02}

\[\theta = - 0.053 + 0.0292K - 5.5 \times 10^{- 4}K^{2} + 4.3 \times 10^{- 6}K^{3}\]

όπου θ η ογκομετρική περιεκτικότητα του εδάφους σε νερό και Κ η
πραγματική συνιστώσα της σχετικής ηλεκτρικής διαπερατότητας.
